\chapter*{Glossary and Acronyms}
\addcontentsline{toc}{chapter}{Glossary and Acronyms}

\noindent This list explains the acronyms and the specialist terms used in this dissertation. Terms
marked \textit{(project term)} have a specific meaning fixed by this project's rules rather than a
general one.

\begin{longtable}{P{0.20\textwidth}P{0.72\textwidth}}
\toprule
\textbf{Term} & \textbf{Meaning} \\
\midrule
\endhead
Agent & A software component given one named job in a workflow, such as collecting evidence or checking a claim. \\
AI & Artificial intelligence. \\
API & Application programming interface: a fixed way for one program to ask another for data. \\
C0--C3 & The four planned comparison conditions: manual work, a deterministic baseline, the multi-agent workflow with verification, and that workflow plus human review. \\
Checksum & A short fingerprint calculated from a file's contents. If one byte changes, the fingerprint changes, so a reader can tell whether a saved file has been altered. Also called a content hash. \\
CLI & Command-line interface: text commands typed into a terminal rather than clicked in a browser. \\
Claim & \textit{(project term)} A structured statement about one company, one measure, one value and one reporting period, which must be checked before it may appear in a report. \\
Claim admission & \textit{(project term)} The decision to let a checked claim enter the report. Refusal of claim admission withholds the claim and records the reason. \\
Contradicted & \textit{(project term)} A verification outcome: eligible evidence disagrees with the proposed value, so the claim is withheld and the disagreement is kept on record. \\
Coverage & See \textit{statement coverage}. \\
CSV, JSON, XLSX & Common file formats for tabular or structured data. \\
D0--D2 & The three dataset tiers: D0 synthetic development cases, D1 an authorised pilot, and D2 a sealed final out-of-sample set. Only D0 produced observations. \\
Deterministic & Producing the same output from the same input every time, with no element of chance. \\
Design science & A research approach in which a working artefact is built to address a stated problem, and both the artefact and the lessons from building it are treated as the contribution. \\
Docker, Docker Compose & Tools that package software with its dependencies into containers so it runs the same way on different machines. \\
Egress & Outbound network traffic leaving a system. \\
Evidence item & \textit{(project term)} A saved passage from a saved source copy, recorded with its exact location, so a later reader can find the same words again. \\
Fixture & A fixed, made-up test input used to check that software behaves as specified. \\
Held & \textit{(project term)} The D0 fixture label for refusal or non-admission: the claim is not supported, is out of date or comes from untrusted content, so it is withheld from the report with its reason recorded. It is not a human hold. \\
Refusal & \textit{(project term)} Withholding claim admission when the saved passage does not support the statement, and recording that outcome instead of rewriting the claim. \\
HTTPS, TLS & The encrypted protocol and the encryption layer used for secure web connections. \\
LLM & Large language model: a model trained to generate text, such as those behind current chat assistants. \\
Loopback & A network address reachable only from the same machine, so the service is not exposed to the network. \\
Manifest & A list that binds outputs to their sources and fingerprints, so an export can be checked afterwards. \\
Nulls, missing states & \textit{(project term)} The eight narrative states are blank, explicit zero, none stated, not applicable, not reported, not found publicly, source unavailable and invalid. The implemented type system also records observed, filing-not-due, dormant and not-required. \\
Prompt injection & Hidden instructions placed in a web page or document that try to make a language model ignore its own rules. \\
Provenance & The recorded origin and handling history of a piece of data: where it came from, when it was saved and what was done to it. \\
RAG & Retrieval-augmented generation: retrieving documents first and then generating text from them, rather than relying on the model's memory. \\
RQ & Research question. \\
Schema & A declared structure that data must match, used to reject malformed output automatically. \\
Snapshot, freeze & A dated, unchanging copy of data or results, taken so later changes cannot alter what was reported. \\
SQLite & A small database stored as a single local file. \\
Statement coverage & The share of a program's lines that were executed at least once by its tests. It shows what was exercised, not that the behaviour is correct. \\
\texttt{store=False} & A setting asking an external model provider not to save the request as application state. It does not remove separate abuse-monitoring retention. \\
Supported & \textit{(project term)} A verification outcome: eligible evidence matches the proposed value, period and units, so the claim may enter the report. \\
Synthetic data & Made-up data written for testing. It contains no real company or personal information, and results from it describe the software, not the world. \\
URL & Uniform resource locator: a web address. \\
Verification & \textit{(project term)} A separate check, carried out after a claim has been drafted, that compares the claim against its evidence and records an outcome without rewriting the claim. \\
\bottomrule
\end{longtable}
